\section{Sard's Theorem}

\subsection{Sets of measure zero}

\bdefn

    Recall that a subset $S\subseteq\bR^n$ is said to be of {\emphcolor measure zero} if for every $\epsilon>0$, $S$ can be covered by
    a countable collection of open rectangles where the sum of their volume is less than $\epsilon$.
    Equivalently, $S$ literally has zero measure with respect to the Lebesgue measure on $\bR^n$.

\edefn

\blemm

    Suppose that $A\subseteq\bR^n$ is measurable (including Borel subsets), such that for every $c\in\bR$,
    $A$'s intersection with $\set c\times\bR^{n-1}$ has $(n-1)$-dimension measure zero.
    Then $A$ has measure zero.

\elemm

\bproof

    $A$ is measurable and so its indicator function $\chi_A$ is integrable.
    By Fubini's theorem,
    $$ \mu(A) = \int_{\bR^n}\chi_A\,d(x^1,\dots,x^n) = \int_\bR\int_{\bR^{n-1}}\chi_A(c,x^2,\dots,x^n)\,d(x^2,\dots,x^n)\,dx^1 , $$
    where $\mu$ is the Lebesgue measure.
    Now the integral of $\chi_A(c,\bar x)$ is the measure of the intersection of $A$ with $\set c\times\bR^{n-1}$, which we assumed to be zero.
    Thus the integral is zero and $A$ has measure zero.
    \qed

\eproof

\bprop

    Let $A$ be open or closed in $\bR^{n-1}$ or $\bH^{n-1}$, and $f\colon A\to\bR$ continuous.
    Then $f$'s graph has measure zero in $\bR^n$.

\eprop

\bproof

    Note that $f$'s graph is closed: it is the preimage of $0$ under the map $g\colon\bR^n\to\bR$ mapping $(\bar x,a)$ to $f(\bar x)-a$.
    For $c\in\bR$, $f$'s graph's intersection with $\set c\times\bR^{n-1}$ is the graph of the function mapping $x\in\bR^{n-2}$ to $f(c,x)$, i.e.\ the graph
    of the restriction of $f$ to $\bR^{n-2}$.
    We can then induct, showing that for every $c\in\bR$ these have measure zero, and then appeal to the above lemma.
    \qed

\eproof

\bcoro

    Every proper affine subspace of $\bR^n$ has measure zero.

\ecoro

\bproof

    If $S$ is an affine subspace of $\bR^n$ of dimension $n-1$, it is the graph of an affine function $\bR^{n-1}\to\bR$.
    By the above proposition, it has measure zero.
    \qed

\eproof

\bprop

    Suppose $A\subseteq\bR^n$ has measure zero and $F\colon A\to\bR^n$ is smooth.
    Then $F(A)$ has measure zero.

\eprop

\bproof

    For $p\in A$, $F$ extends to an open ball around $p$.
    By potentially shrinking the ball, we can extend $F$ to its closure, a compact set $\bar U$.
    $A$ can be covered by countably many $A\cap\bar U$s, so it is sufficient to show that $F(A\cap\bar U)$ has measure zero.

    Since $\bar U$ is compact, $\norm{JF(x)}$ is bound by some constant $C$ on it (this is true for all continuous real-valued functions
    on compact sets).
    Then we see that $\norm{F(x)-F(y)}\leq C\norm{x-y}$ for all $x,y\in\bar U$.
    Since $A$ has measure zero, for any $\epsilon>0$ we can cover $A\cap\bar U$ with balls $\set{B_j}$ of total measure $<\epsilon$.
    By Lipschitz continuity, $F(B_j)$'s measure is bound by $C$ times $B_j$'s, so that
    $$ \mu F(A\cap\bar U) \leq \sum_j \mu F(B_j) \leq \sum_jC\mu(B_j) < C\epsilon . $$
    Since $\epsilon$ is arbitrary, we have the desired result.
    \qed

\eproof

\bnote

    In the above proof we used the fact that smooth functions from a compact set are Lipschitz continuous.
    Indeed, if $F\colon\bar U\to\bR^n$ is smooth (even $C^1$), then let $\norm{JF(x)}\leq C$ for all $x\in C$, then
    $$ F(y) - F(x) = \int_0^1\frac d{dt}F(x+t(y-x))\,dt $$
    by the fundamental theorem of calculus.
    Now,
    $$ \frac d{dt}F(x+t(y-x)) = (y-x)JF(x+t(y-x)) $$
    so we have that
    $$ \norm{F(y)-F(x)} = \norm{\int_0^1 (y-x)JF(x+t(y-x))\,dt} \leq \int_0^1\norm{y-x}\norm{JF(x+t(y-x))}\,dt . $$
    We can bound the Jacobian by $C$, giving us
    $$ \leq \norm{y-x} \int_0^1 C = C\norm{y-x} $$
    as desired.

\enote

\bdefn

    If $M$ is a smooth $n$-manifold, with or without boundary, a subset $A\subseteq M$ has {\emphcolor measure zero} if for every smooth chart
    $(U,\phi)$ for $M$, $\phi(A\cap U)\subseteq\bR^n$ has measure zero.

\edefn

\blemm

    Let $(U_\alpha,\phi_\alpha)$ be a collection of smooth charts covering $M$.
    If $A\subseteq M$ is a subset where $\phi(A\cap U_\alpha)$ has measure zero for all $\alpha$, then $A$ has measure zero.

\elemm

\bproof

    Let $(V,\psi)$ be an arbitrary smooth chart, we need to show $\psi(A\cap V)$ has measure zero.
    For each $\alpha$, consider
    $$ \psi(A\cap V\cap U_\alpha) = (\psi\circ\phi_\alpha)^{-1}\circ\phi_\alpha(A\cap V\cap U_\alpha) , $$
    which has measure zero by the above proposition.
    Since $\psi(A\cap V)$ is covered by countably many of these, it has measure zero.
    \qed

\eproof

\bprop

    Let $A\subseteq M$ be a subset with measure zero.
    Then $M-A$ is dense in $M$.

\eprop

\bproof

    If not, then there is an open subset $U$ contained in $A$
    Let $(V,\phi)$ be a chart contained in $U\subseteq V$, then $\phi(A\cap V)=\phi(V)$, which is nonempty open in $\bR^n$.
    But open subsets of $\bR^n$ cannot have measure zero (they contain an open ball/rectangle which by definition has positive measure),
    a contradiction.
    \qed

\eproof

\bthrm

    Let $F\colon M\to N$ be smooth, and $A\subseteq M$ of measure zero.
    Then $F(A)\subseteq N$ has measure zero.

\ethrm

\bproof

    Let $(U_\alpha,\phi_\alpha)$ cover $M$, and let $(V,\psi)$ be a chart for $N$.
    Then $\psi(F(A)\cap V)$ can be covered by $\psi\circ F\circ\phi_\alpha^{-1}(\phi_\alpha(A\cap U_\alpha\cap F^{-1}(V)))$.
    Now $\psi\circ F\circ\phi_\alpha^{-1}$ is a smooth Euclidean map, and $\phi_\alpha(A\cap U_\alpha\cap F^{-1}(V))$ has measure zero by assumption,
    so these coverings have measure zero.
    Thus so does $\psi(F(A)\cap V)$ as desired.
    \qed

\eproof

\bdefn

    Let $F\colon M\to N$ be smooth.
    A point $p\in M$ is a {\emphcolor regular point} of $F$ if $F$ has rank $\dim N$ at $p$, otherwise it is a {\emphcolor critical point}.
    The image of a regular point is called a {\emphcolor regular value}, and the image of a critical point is a {\emphcolor critical value}.

\edefn

In other words, a point is a regular point iff $dF_p\colon T_pM\to T_{Fp}N$ is a surjection.
Notice that if $\dim M<\dim N$, every point is critical.

\bthrm[title=Sard's theorem]

    Let $F\colon M\to N$ be smooth map.
    Then the set of critical values has measure zero in $N$.

\ethrm

\bproof

    Let $m=\dim M,n=\dim N$, we induct on $m$.
    For $m=0$ it is clear: either $n=0$ for which $F$ has no critical points, or $n>0$ and all points are critical but $M$ is countable and
    thus $F(M)\subseteq N$ has measure zero.

    By covering $F$ with countably many smooth charts, we can reduce this to the case that $F$ is a smooth map from an open subset $U$ of $\bR^m$
    or $\bH^m$ to $\bR^n$.
    Let $C\subseteq U$ be the set of critical points of $F$.
    Define
    $$ C_k = \set{x\in C}[\hbox{For $1\leq i\leq k$, all of the $i$th partial derivatives of $F$ vanish at $x$}] . $$
    Then $C_k$ is decreasing, $C\supseteq C_1\supseteq C_2\supseteq\cdots$, and by continuity they are all closed in $U$.

    Now we prove in three parts that $F(C)$ has measure zero.
    \benum
        \item Firstly, $F(C-C_1)$ has measure zero.
        Because $C_1$ is closed in $U$, we can take $U=U-C_1$ and assume $C_1=\emptyset$.
        Let $a\in C$, so we assumed that $F$'s first partial derivatives are not all $0$ at $a$.
        Without loss of generality, $\ipdvof F^1{x^1}(a)\neq0$.
        So we can define new smooth coordinates $(u,v)=(u,v^2,\dots,v^m)$ in a neighborhood of $a$, where $u=F^1$ and $v^i=x^i$.
        This is because the Jacobian of these coordinates is invertible at $a$ and thus locally a diffeomorphism.
        Let $V_a$ be the neighborhood of $a$ using these coordinates, and potentially shrinking it, we can assume $\cl V_a$ is compact in $U$.

        In these coordinates, $F$ has representation
        $$ F(u,v^2,\dots,v^m) = (u,F^2(u,v),\dots,F^n(u,v)) , $$
        so its Jacobian is
        $$ JF(u,v) = \pmatrix{1&0\cr *&\pdvof{F^i}{v^j}} . $$
        Thus $C\cap\cl V_a$ consists of points the $(n-1)\times(m-1)$ matrix $\pdvof{F^i}{v^j}$ has rank $<n-1$.

        We wish to show that $F(C\cap\cl V_a)$ has measure zero in $N$.
        This set is compact, and thus measurable, so it is sufficient to show that its intersection with the hyperplane $y^1=c$ has $n-1$-dimensional
        measure zero.
        For $c\in\bR$, define $B_c=\set v[(c,v)\in\cl V_a]\subseteq\bR^{m-1}$, and define $F_c\colon B_c\to\bR^{n-1}$ by
        $$ F_c(v) = (F^2(c,v),\dots,F^n(c,v)) . $$
        Because $F(c,v)=(c,F_c(v))$, the critical values of $F|_{\cl V_a}$ which lie in the hyperplane $y^1=c$ are precisely the points of the form $(c,w)$
        where $w$ is a critical value of $F_c$.
        By induction, these have $n-1$-dimensional measure zero, so $F(C\cap\cl V_a)$ has measure zero.

        We can cover $U$ with countably many sets of the form $\cl V_a$, and thus $C$ by countably many sets of the form $C\cap\cl V_a$.
        Thus $F(C)$ is the countable union of sets of measure zero, and thus has measure zero.

        \item For the second step, we claim $F(C_k-C_{k+1})$ has measure zero for each $k$.
        As before, we can take $U=U-C_{k+1}$ and assume that every point in $C_k$ has a $k+1$th partial derivative which does not vanish.
        Let $a\in C_k$, and denote by $y\colon U\to\bR$ some $k$th partial derivative of $F$ which has a nonvanishing partial derivative at $a$.
        Since $a\in C_k$ we have that $y(a)=0$.

        Then $a$ is a regular point of $y$, so it has a neighborhood $V_a$ consisting of regular points of $y$.
        Let $Y$ be the zero set of $y$ in $V_a$, i.e.\ $Y=y^{-1}0$.
        Since $V_a$ are regular points of $y$, by definition $Y$ is a regular level set, and thus an embedded submanifold of $V_a$ of codimension $1$.
        Since $V_a$ is open, $Y$ is an embedded submanifold of $M$ of dimension $m-1$.

        For any $p\in C_k\cap V_a$, $dF_p$ is not surjective (by assumption, $p$ is critical), so then neither is $d(F|_Y)_p=(dF_p)|_{T_pY}$.
        So $F(C_k\cap V_a)$ is contained in the critical values of $F|_Y\colon Y\to\bR^n$, which by induction has measure zero.
        Since $U$ can be covered by countably many $V_a$, the result follows.

        \item Now, notice that the intersection of all $C-C_1$ and $C_k-C_{k+1}$ need not be all of $C$: we may have points in $C$ all of whose partial
        derivatives vanish.
        So for our final step, we claim that for $k>m/n+1$, $F(C_k)$ has measure zero.

        For $a\in U$, let $E\subseteq U$ be a closed cube containing $a$.
        Since $U$ can be covered by countably many such cubes, we will show $F(C_k\cap E)$ has measure zero for all such closed cubes.

        Since $E$ is compact, we can find a value $A>0$ which bounds the absolute values of all of $F$'s $k+1$th partial derivatives in $E$.
        Let $R$ be the side length of $E$, and let $K$ be some integer to be chosen.
        Then we can subdivide $E$ into $K^m$ subcubes of side length $R/K$, denoted by $(E_1,\dots,E_{K^m})$.
        If $E_i$ is a cube and there is an $a_i\in E_i\cap C_k$, then for all $x\in E_i$ we have
        $$ \norm{F(x) - F(a_i)} \leq A'\norm{x-a_i}^{k+1} $$
        for some value $A'$ dependent only on $A,k,m$.
        This means that $F(E_i)$ is contained in a ball of radius $A'(R/K)^{k+1}$, so of measure $A''(R/K)^{nk+n}$.
        Then $F(C_k\cap E)$ is contained in the union of $K^m$ balls of this radius.
        So its measure is bound by
        $$ K^m\cdot A''(R/K)^{nk+n} = A''R^{nk+n}K^{nk+n-m} . $$
        By assumption $nk+n-m<1$, so this can be made arbitrarily small by taking large enough $K$, so $F(C_k\cap E)$ has measure zero.
    \eenum

    Now, we just note that
    $$ C = (C-C_1) \cup \bigcup_k (C_k-C_{k+1}) \cup \bigcup_{k>m/n+1} C_k , $$
    a countable union of sets of measure zero, and thus $C$ has measure zero.
    \qed

\eproof

\bnote

    In the above proof we had a point where we had a bound like
    $$ \norm{F(x) - F(a)} \leq A\norm{x-a}^{k+1} . $$
    We obtain this bound by considering each of $F$'s coordinate's Taylor expansions.

    If $f\colon W\subseteq\bR^n\to\bR$ is smooth and defined in some convex subset $W$, then take $f$'s $k$th order Taylor expansion $P_k(x)$ around $a\in W$.
    Its remainder can be bound by
    $$ \abs{f(x) - P_k(x)} \leq \frac{n^{k+1}M}{(k+1)!}\norm{x-a}^{k+1} , $$
    as well as
    $$ \abs{P_k(x) - f(a)} \leq \frac{kn^kM}{k!}\norm{x-a}^k $$
    (Just take a good look at the remainder and polynomial expressions, here $M$ is a bound on the partial derivatives.)
    Then we have
    $$ \abs{f(x) - f(a)} \leq \abs{f(x) - P_k(x)} + \abs{P_k(x) - f(a)} \leq A\norm{x-a}^{k+1} $$
    for some absolute constant $A$ dependent only on $n,k$, and $M$ (a bound on the partial derivatives).
    We then get that
    $$ \norm{F(x) - F(a)} = \sqrt{\sum_i\abs{F^i(x)-F^i(a)}^2} \leq nA\norm{x-a}^{k+1} . $$

\enote

\subsection{Transversality}

\bdefn

    Let $F\colon M\to N$ be a smooth map and $S\subseteq N$ an embedded submanifold.
    We say $F$ is {\emphcolor transverse} to $S$ if for every $x\in F^{-1}(S)$, $T_{Fx}S$ and $dF_x(T_xM)$ span
    $T_{Fx}N$.

    Two embedded submanifolds $S,S'\subseteq M$ {\emphcolor intersect transversally} if for each $p\in S\cap S'$,
    $T_pS$ and $T_pS'$ span $T_pM$.

\edefn

Note that $S,S'$ intersect transversally iff the inclusion of one in $M$ is transverse to the other.

Furthermore if $F$ is a submersion, then $dF_x$ is surjective and thus $F$ is transverse to every embedded submanifold.

\bthrm

    Let $M,N$ be smooth manifolds, $S\subseteq N$ an embedded submanifold.
    \benum
        \item If $F\colon M\to N$ is transverse to $S$ then $F^{-1}S$ is an embedded submanifold of $M$ whose codimension
        in $M$ is equal to $S$'s codimension in $N$.
        \item If $S'\subseteq N$ is an embedded submanifold that intersects $S$ transversally, $S\cap S'$ is an embedded
        submanifold of $N$ whose codimension is equal to the sum of the codimensions of $S$ and $S'$.
    \eenum

\ethrm

\bproof

    The second point follows from the first, by taking $F$ to be the inclusion $S'\inj N$.
    Then $F^{-1}S=S\cap S'$ is an embedded submanifold of $S'$ of codimension $\codim_NS$ by the first point.
    The composition of embeddings is an embedding, so $F^{-1}S$ is embedded in $N$, and its codimension is
    $$ \multline{
        \codim_N(S\cap S') = \dim N - \dim S' + \dim S' - \dim(S\cap S') = \codim_NS' + \codim_{S'}(S\cap S')\cr
        = \codim_NS' + \codim_NS' .
    } $$

    Let $n=\dim N$ and $k=\codim_NS$.
    For $x\in F^{-1}(S)$, we find a neighborhood $U$ of $F(x)$ in $M$ and a local defining function $\phi\colon U\to\bR^k$
    such that $S\cap U=\phi^{-1}0$.
    If we can show that $0$ is a regular value of $\phi\circ F$, because $F^{-1}(S)\cap F^{-1}(U)$ is the zero set
    of $\phi\circ F|_{F^{-1}U}$.

    For $z\in T_0\bR^k$ and $p\in(\phi\circ F)^{-1}0$, $0$ being a regular value of $\phi$ means that there is a vector
    $y\in T_{Fp}N$ such that $d\phi_{Fp}y=z$.
    $F$ being transverse to $S$ means that $y=y_0+dF_pv$ for some $y_0\in T_{Fp}S$ and $v\in T_pM$.
    Because $\phi$ is constant on $S\cap U$, it follows that $d\phi_{Fp}y_0=0$ so
    $$ z = d\phi_{Fp}y = d\phi_{Fp}y_0 + d\phi_{Fp}dF_pv = d(\phi\circ F)_pv . $$
    So $\phi\circ F$ is indeed a submersion.
    \qed

\eproof

Note that this proof showed that $x\in F^{-1}S$ is a regular point of $F$.

\bprop

    If $F\colon M\to N$ is transverse to $S\subseteq N$, then the points of $F^{-1}S$ are regular points of $F$.

\eprop

Since submersions are transverse to every embedded submanifold, the following corollary is immediate.

\bcoro

    If $S\subseteq N$ is an embedded submanifold, and $F\colon M\to N$ is smooth, then $F^{-1}S$ is an embedded submanifold
    of $M$ with
    $$ \codim_MF^{-1}S = \codim_NS . $$

\ecoro

\bdefn

    Suppose $N,M,S$ are smooth manifolds, and for each $s\in S$ we are given a map $F_s\colon N\to M$.
    The collection $\set{F_s}_{s\in S}$ is called a {\emphcolor smooth family of maps} if $F\colon M\times S\to N$
    defined by $F(m,s)=F_s(m)$ is smooth.

\edefn

The notion of a smooth family of maps is a sort of higher-dimensional analogue of homotopy.

\bprop

    If $\set{F_s}_{s\in S}$ is a smooth family of maps where $S$ is connected, then for every $s_1,s_2\in S$, $F_{s_1},F_{s_2}$
    are homotopic.

\eprop

\bproof

    Since $S$ is connected, let $\gamma\colon I\to S$ be a path from $s_1$ to $s_2$.
    Then $H\colon M\times I\to N$ given by $H(m,t)=F_{\gamma(t)}(m)$ is a homotopy from $F_{s_1}$ to $F_{s_2}$.
    \qed

\eproof

\blemm

    If $F\colon N\to M$ is transverse to $X\subseteq M$ embedded, let $W=F^{-1}X$.
    Then for every $p\in W$, $dF_p(T_pW)=T_{Fp}X$ and $T_pW=(dF_p)^{-1}(T_{Fp}X)$.
    Thus if $X,X'\subseteq M$ intersect transversally then $T_p(X\cap X')=T_pX\cap T_pX'$ for all $p\in X\cap X'$.

\elemm

\bproof

    The second point is due to $X\cap X'$ being the preimage of $X'$ under the inclusion $X\inj M$.
    Then $(d\iota_p)^{-1}(T_pX')=T_pX\cap T_pX'$.

    The first point is because $F$ maps $W$ into $X$ so $dF_p(T_pW)\subseteq T_{Fp}X$ for all $p\in W$.
    Now by rank-nullity,
    $$ \dim dF_p(T_pW) = \dim T_pW - \dim\ker dF_p = \dim W - \dim\ker dF_p . $$
    Now recall that $p$ is a regular point of $F$, so $dF_p$ is surjective and thus
    $$ \dim\ker dF_p = \dim T_pN - \dim T_{Fp}M = \dim N - \dim M , $$
    so
    $$ \dim dF_p(T_pW) = \dim W - \dim N + \dim M = \dim M - \codim_NW = \dim X , $$
    since $\codim_MX=\codim_NW$, and thus considering dimensions $dF_p(T_pW)=T_{Fp}X$.

    And so $T_pW\subseteq(dF_p)^{-1}(T_{Fp}X)$, as well, and considering dimensions they are equal.
    Indeed, $\dim T_pW=\dim W$ and in general if $T\colon V\to U$ is a surjective linear map, and $U'\subseteq U$ a subspace,
    then
    $$ \dim T^{-1}U' = \dim V - \dim U' + \dim T(U') . $$
    This is because $T|_{T^{-1}U'}\colon T^{-1}U'\to U'$ has kernel $\ker T$, and thus by rank-nullity
    $$ \dim T^{-1}U' = \dim\ker T + \dim U' = \dim V - \dim U + \dim U' . $$
    $dF_p$ is surjective onto its image $dF_p(T_pM)$, so
    $$ \dim(dF_p)^{-1}T_{Fp}X = \dim T_pN - \dim dF_p(T_pN) + \dim T_{Fp}X = \dim N - \dim dF_p(T_pN) + \dim X . $$
    Since $p\in F^{-1}X$, it is a regular point of $F$, and so $dF_p(T_pN)=T_pM$.
    Thus this is equal to
    $$ = \dim N - \dim M + \dim X = \dim W . $$
    \qed

\eproof

\bdefn

    A subset of a manifold is said to contain {\emphcolor almost every element} if its complement has measure zero.

\edefn

\bthrm[title=Parametric transversality]

    Suppose $N,M$ are smooth manifolds, $X\subseteq M$ an embedded submanifold, and $\set{F_s}_{s\in S}$
    a smooth family of maps from $N$ to $M$.
    If $F\colon N\times S\to M$ is transverse to $X$, then for almost every $s\in S$, $F_s$ is transverse to $X$.

\ethrm

\bproof

    By assumption $W=F^{-1}(X)$ is an embedded submanifold of $N\times S$.
    Let $\pi\colon N\times S\to S$ be projection onto the second component.
    We will show that if $s\in S$ is a regular value of $\pi|_W$, then $F_s$ is transverse to $X$.
    By Sard's theorem almost every $s\in S$ is a regular value, so the theorem will be proven.

    If $s\in S$ is a regular value of $\pi|_W$, then let $p\in F_s^{-1}(X)$ be arbitrary, and set $q=F_s(p)$.
    We want to show that $T_qM=T_qX+d(F_s)(T_pN)$.
    By our hypothesis on $F$,
    $$ T_qM = T_qX + dF(T_{(p,s)}(N\times S)) , $$
    and since $s$ is a regular value of $\pi|_W$ and $(p,s)\in W$, we have
    $$ T_sS = d\pi(T_{(p,s)}W) . $$
    The above lemma showed that
    $$ dF_{(p,s)}(T_{(p,s)}W) = T_qX . $$

    We want to show that $F_s$ is transverse to $X$, i.e.
    $$ T_qM = T_qX + dF_s(T_pN) . $$
    So let $w\in T_qM$ arbitrary, we want $v\in T_qX$ and $y\in T_pN$ such that $w=v+dF_s(y)$.
    Since $F$ is transverse to $X$, there is $v_1\in T_qX$ and $(y_1,z_1)\in T_pN\times T_sS\cong T_{(p,s)}(N\times S)$ such that
    $$ w = v_1 + dF(y_1,z_1) . $$
    And then since $T_sS=d\pi(T_{(p,s)}W)$ there is a $(y_2,z_2)\in T_{(p,s)}W$ such that $d\pi(y_2,z_2)=z_1$.
    Since $\pi$ is a projection, this means $z_2=z_1$.
    By linearity,
    $$ dF(y_1,z_1) = dF(y_2,z_1) + dF(y_1-y_2,0) . $$

    Now,
    $$ dF(y_2,z_1) = dF(y_2,z_2) \in dF(T_{(p,s)}W) = T_qX . $$
    On the other hand let $\iota_s\colon N\to N\times S$ be the slice map $\iota_s(n)=(n,s)$ so $F_s=F\circ\iota_s$ and
    $d\iota_s(y_1-y_2)=(y_1-y_2,0)$ so
    $$ dF(y_1-y_2,0) = dF\circ d\iota_s(y_1-y_2) = dF_s(y_1-y_2) . $$
    And so all in all we have
    $$ w = (v_1 + dF(y_2,z_1)) + dF_s(y_1-y_2) , $$
    where $v_1,dF(y_2,z_1)\in T_qX$ and $y_1-y_2\in T_pN$, as desired.
    \qed

\eproof

\bthrm[title=Transversality homotopy]

    If $F\colon N\to M$ is smooth and $X\subseteq M$ is embedded, then $F$ is homotopic to a smooth map $G\colon N\to M$
    which is transverse to $X$.

\ethrm

\bproof

    The idea of the proof is to define a smooth map $H\colon N\times S\to M$ which is transverse to $X$, where $S$ is connected, and $H_s=F$ for
    some $s\in S$.
    Then by parametric transversality there is a $t\in S$ such that $H_t$ is transverse to $X$.
    And $F=H_s$ and $G=H_t$ are homotopic, as desired.

    We omit the details.
    \qed

\eproof

